% Options for packages loaded elsewhere
\PassOptionsToPackage{unicode}{hyperref}
\PassOptionsToPackage{hyphens}{url}
\documentclass[
]{article}
\usepackage{xcolor}
\usepackage[margin=1in]{geometry}
\usepackage{amsmath,amssymb}
\setcounter{secnumdepth}{-\maxdimen} % remove section numbering
\usepackage{iftex}
\ifPDFTeX
  \usepackage[T1]{fontenc}
  \usepackage[utf8]{inputenc}
  \usepackage{textcomp} % provide euro and other symbols
\else % if luatex or xetex
  \usepackage{unicode-math} % this also loads fontspec
  \defaultfontfeatures{Scale=MatchLowercase}
  \defaultfontfeatures[\rmfamily]{Ligatures=TeX,Scale=1}
\fi
\usepackage{lmodern}
\ifPDFTeX\else
  % xetex/luatex font selection
\fi
% Use upquote if available, for straight quotes in verbatim environments
\IfFileExists{upquote.sty}{\usepackage{upquote}}{}
\IfFileExists{microtype.sty}{% use microtype if available
  \usepackage[]{microtype}
  \UseMicrotypeSet[protrusion]{basicmath} % disable protrusion for tt fonts
}{}
\makeatletter
\@ifundefined{KOMAClassName}{% if non-KOMA class
  \IfFileExists{parskip.sty}{%
    \usepackage{parskip}
  }{% else
    \setlength{\parindent}{0pt}
    \setlength{\parskip}{6pt plus 2pt minus 1pt}}
}{% if KOMA class
  \KOMAoptions{parskip=half}}
\makeatother
\usepackage{color}
\usepackage{fancyvrb}
\newcommand{\VerbBar}{|}
\newcommand{\VERB}{\Verb[commandchars=\\\{\}]}
\DefineVerbatimEnvironment{Highlighting}{Verbatim}{commandchars=\\\{\}}
% Add ',fontsize=\small' for more characters per line
\usepackage{framed}
\definecolor{shadecolor}{RGB}{248,248,248}
\newenvironment{Shaded}{\begin{snugshade}}{\end{snugshade}}
\newcommand{\AlertTok}[1]{\textcolor[rgb]{0.94,0.16,0.16}{#1}}
\newcommand{\AnnotationTok}[1]{\textcolor[rgb]{0.56,0.35,0.01}{\textbf{\textit{#1}}}}
\newcommand{\AttributeTok}[1]{\textcolor[rgb]{0.13,0.29,0.53}{#1}}
\newcommand{\BaseNTok}[1]{\textcolor[rgb]{0.00,0.00,0.81}{#1}}
\newcommand{\BuiltInTok}[1]{#1}
\newcommand{\CharTok}[1]{\textcolor[rgb]{0.31,0.60,0.02}{#1}}
\newcommand{\CommentTok}[1]{\textcolor[rgb]{0.56,0.35,0.01}{\textit{#1}}}
\newcommand{\CommentVarTok}[1]{\textcolor[rgb]{0.56,0.35,0.01}{\textbf{\textit{#1}}}}
\newcommand{\ConstantTok}[1]{\textcolor[rgb]{0.56,0.35,0.01}{#1}}
\newcommand{\ControlFlowTok}[1]{\textcolor[rgb]{0.13,0.29,0.53}{\textbf{#1}}}
\newcommand{\DataTypeTok}[1]{\textcolor[rgb]{0.13,0.29,0.53}{#1}}
\newcommand{\DecValTok}[1]{\textcolor[rgb]{0.00,0.00,0.81}{#1}}
\newcommand{\DocumentationTok}[1]{\textcolor[rgb]{0.56,0.35,0.01}{\textbf{\textit{#1}}}}
\newcommand{\ErrorTok}[1]{\textcolor[rgb]{0.64,0.00,0.00}{\textbf{#1}}}
\newcommand{\ExtensionTok}[1]{#1}
\newcommand{\FloatTok}[1]{\textcolor[rgb]{0.00,0.00,0.81}{#1}}
\newcommand{\FunctionTok}[1]{\textcolor[rgb]{0.13,0.29,0.53}{\textbf{#1}}}
\newcommand{\ImportTok}[1]{#1}
\newcommand{\InformationTok}[1]{\textcolor[rgb]{0.56,0.35,0.01}{\textbf{\textit{#1}}}}
\newcommand{\KeywordTok}[1]{\textcolor[rgb]{0.13,0.29,0.53}{\textbf{#1}}}
\newcommand{\NormalTok}[1]{#1}
\newcommand{\OperatorTok}[1]{\textcolor[rgb]{0.81,0.36,0.00}{\textbf{#1}}}
\newcommand{\OtherTok}[1]{\textcolor[rgb]{0.56,0.35,0.01}{#1}}
\newcommand{\PreprocessorTok}[1]{\textcolor[rgb]{0.56,0.35,0.01}{\textit{#1}}}
\newcommand{\RegionMarkerTok}[1]{#1}
\newcommand{\SpecialCharTok}[1]{\textcolor[rgb]{0.81,0.36,0.00}{\textbf{#1}}}
\newcommand{\SpecialStringTok}[1]{\textcolor[rgb]{0.31,0.60,0.02}{#1}}
\newcommand{\StringTok}[1]{\textcolor[rgb]{0.31,0.60,0.02}{#1}}
\newcommand{\VariableTok}[1]{\textcolor[rgb]{0.00,0.00,0.00}{#1}}
\newcommand{\VerbatimStringTok}[1]{\textcolor[rgb]{0.31,0.60,0.02}{#1}}
\newcommand{\WarningTok}[1]{\textcolor[rgb]{0.56,0.35,0.01}{\textbf{\textit{#1}}}}
\usepackage{graphicx}
\makeatletter
\newsavebox\pandoc@box
\newcommand*\pandocbounded[1]{% scales image to fit in text height/width
  \sbox\pandoc@box{#1}%
  \Gscale@div\@tempa{\textheight}{\dimexpr\ht\pandoc@box+\dp\pandoc@box\relax}%
  \Gscale@div\@tempb{\linewidth}{\wd\pandoc@box}%
  \ifdim\@tempb\p@<\@tempa\p@\let\@tempa\@tempb\fi% select the smaller of both
  \ifdim\@tempa\p@<\p@\scalebox{\@tempa}{\usebox\pandoc@box}%
  \else\usebox{\pandoc@box}%
  \fi%
}
% Set default figure placement to htbp
\def\fps@figure{htbp}
\makeatother
\setlength{\emergencystretch}{3em} % prevent overfull lines
\providecommand{\tightlist}{%
  \setlength{\itemsep}{0pt}\setlength{\parskip}{0pt}}
\usepackage{bookmark}
\IfFileExists{xurl.sty}{\usepackage{xurl}}{} % add URL line breaks if available
\urlstyle{same}
\hypersetup{
  pdftitle={HW6},
  pdfauthor={Nathan Krieger},
  hidelinks,
  pdfcreator={LaTeX via pandoc}}

\title{HW6}
\author{Nathan Krieger}
\date{2026-03-06}

\begin{document}
\maketitle

\section{Problem 1: Concept Review}\label{problem-1-concept-review}

\subsection{(a)}\label{a}

Best subset selection will produce a collection of p models M1, M2, . .
. , Mp. These represent the `best' model of each size (where `best' here
is defined as the model with the smallest RSS).

\subsubsection{(i) True or False}\label{i-true-or-false}

Is it true that the model identified as Mk+1 must contain a subset of
the predictors found in Mk? In other words, is it true that if M1 : Y
\textasciitilde{} X1, then M2 must also contain X1. And if M2 contains
X1 and X2, then M3 must also contain X1 and X2?

\subsection{(b)}\label{b}

Same question as part (a) but instead of best subset selection, we now
carry out forward selection.

\subsection{(c) (practice)}\label{c-practice}

Suppose we perform best subset selection, forward selection, and
backward selection on a training data set.

Best subset selection produces a model with k predictors for each k = 1,
. . . , p.~Denote this model by Mk,subset. For example, when k = 5,
M5,subset is the best 5-predictor model according to best subset
selection.

Similarly, forward selection produces a model with k predictors, denoted
by Mk,forward, and backward selection produces a model with k
predictors, denoted by Mk,backward.

For each case below, state which of the three models would be expected
to have the smallest training MSE, and briefly explain why.

\subsubsection{(i) (practice) k = 1}\label{i-practice-k-1}

Solution: For k = 1, subset has the same model as forward stepwise and
therefore the same RSS (and training MSE). Backward stepwise may have a
different model, and a worse RSS.

\subsubsection{\texorpdfstring{(ii) (practice)k \({\epsilon}\) \{2,
\ldots, p -
2\}}{(ii) (practice)k \{\textbackslash epsilon\} \{2, \ldots, p - 2\}}}\label{ii-practicek-epsilon-2-p---2}

\subsubsection{(iii) (practice)k = p - 1}\label{iii-practicek-p---1}

\subsubsection{(iv) (practice)k = p}\label{iv-practicek-p}

\subsection{(d) (practice)}\label{d-practice}

Same setup as part (c). For a given k, which of these three models has
the smallest test MSE? Explain your answer.

\subsubsection{(i) k = 1}\label{i-k-1}

Solution: Subset and forward stepwise have the same model when k = 1, so
they have the same test MSE. Backward selection may produce a different
model, so its test MSE may differ.

\subsubsection{(ii) k = p}\label{ii-k-p}

\#\#\#(iii) k = p - 1

\section{Problem 2: Simulation
Studies}\label{problem-2-simulation-studies}

We have seen that as the number of features used in a model increases,
the training error will necessarily decrease, but the test error may
not. We will now explore this in a simulated data set. Start by
generating a data set with p = 20 features, n = 1000 observations, and
an associated response vector Y generated according to the model:

Y = \$\beta\$0 + \$\beta\$1X1 + \$\beta\$2X2 + . . . \(\beta\)pXp +
\(\epsilon\),

You can use this code to get started:

\begin{Shaded}
\begin{Highlighting}[]
\NormalTok{p }\OtherTok{=} \DecValTok{20}
\NormalTok{n }\OtherTok{=} \DecValTok{1000}
\NormalTok{error }\OtherTok{=} \FunctionTok{rnorm}\NormalTok{(n,}\DecValTok{0}\NormalTok{,}\DecValTok{1}\NormalTok{)}
\NormalTok{beta }\OtherTok{=} \FunctionTok{c}\NormalTok{(}\DecValTok{2}\NormalTok{, }\FunctionTok{rep}\NormalTok{(}\DecValTok{0}\NormalTok{,}\DecValTok{5}\NormalTok{),}\FloatTok{1.2}\NormalTok{,}\DecValTok{3}\NormalTok{,}\SpecialCharTok{{-}}\DecValTok{2}\NormalTok{,}\FloatTok{1.5}\NormalTok{,}\DecValTok{7}\NormalTok{,}\FunctionTok{rep}\NormalTok{(}\DecValTok{0}\NormalTok{,}\DecValTok{5}\NormalTok{),}\DecValTok{3}\NormalTok{,}\DecValTok{4}\NormalTok{,}\DecValTok{1}\NormalTok{,}\DecValTok{9}\NormalTok{,}\DecValTok{2}\NormalTok{)}
\NormalTok{Xmat }\OtherTok{=} \FunctionTok{matrix}\NormalTok{(}\FunctionTok{rnorm}\NormalTok{(n}\SpecialCharTok{*}\DecValTok{20}\NormalTok{),n,}\DecValTok{20}\NormalTok{)}
\NormalTok{X }\OtherTok{=} \FunctionTok{cbind}\NormalTok{(}\FunctionTok{rep}\NormalTok{(}\DecValTok{1}\NormalTok{,n),Xmat)}
\NormalTok{Y }\OtherTok{=}\NormalTok{ X}\SpecialCharTok{\%*\%}\NormalTok{beta }\SpecialCharTok{+}\NormalTok{ error}
\NormalTok{data }\OtherTok{=} \FunctionTok{data.frame}\NormalTok{(Y,Xmat)}
\end{Highlighting}
\end{Shaded}

Note that X\%\emph{\%beta is just doing matrix multiplication: beta0 +
beta1}X1 + beta2*X2 + \ldots.

beta20*X20.

\subsection{(a)}\label{a-1}

Split your data set into a training set containing 100 observations and
a test set containing 900 observations.

\subsection{(b)}\label{b-1}

Perform best subset selection on the training set only, and plot the
training set MSE associated with the best model of each size. Hint: The
training MSE can be computed from the model output. If the fitted model
is stored as fit, you can extract it using mean(fit\$rss).

\subsection{(c)}\label{c}

Plot the test MSE associated with the best model of each size.

\subsection{(d)}\label{d}

For which model size does the test MSE take on its minimum values?
Comment on your results.

\subsection{(e)}\label{e}

How does the model at which the test MSE is minimized compare to the
true model used to generate the data? Comment on how well the estimated
regression coefficients match the true values.

Note: The true model contains only 10 predictors with nonzero
coefficients. The remaining 10 predictors have true coefficient value
equal to zero.

\section{Problem 3: Forward and backward selection
(practice)}\label{problem-3-forward-and-backward-selection-practice}

We will use the College data set in the ISLR2 library to predict the
number of applications (Apps) each university received. Use all
remaining variables as predictors.

\subsection{(a) (practice)}\label{a-practice}

Randomly split the data set so that 90\% of the data belong to the
training set and the remaining 10\% belong to the test set. Implement
forward and backward selection on the training set only. Report the test
MSE for all models (M1, M2, . . . , Mp) based on forward selection.
Which model minimizes test MSE?

Report the test MSE for all models (M1, M2, . . . , Mp) based on
backward selection. Which model minimizes test MSE?

Do both algorithms lead you to the same model based on test MSE?

\subsection{(b) (practice)}\label{b-practice}

Run forward and backward selection on the full dataset. For each
approach, report the best model based on AIC. Do they lead you to the
same model based on AIC?

\subsection{(c) (practice)}\label{c-practice-1}

Which is the final model that you would pick? If both (a) and (b) lead
you to the same model, then this is easy. Otherwise, there is no
absolute right answer. Try to justify your choice using some
ideas/concepts you've learned from this class.

\section{Problem 4: A Puzzling
Problem}\label{problem-4-a-puzzling-problem}

When fitting a linear regression model on a data set, you encounter the
following R output. You notice there is something strange about the
results. Point out what is strange in this output and explain clearly
how this could happen.

Call: lm(formula = y \textasciitilde{} x1 + x2)

\section{Problem 5: Interaction
Terms}\label{problem-5-interaction-terms}

We will use the Credit dataset for this problem. It is part of the
library ISLR2

\subsection{(a)}\label{a-2}

This data set contains a few categorical predictors. As we already
discussed in lecture, these predictors should be stored as factors so
that R can handle them properly. Using the str function, check that all
the qualitative predictors in our dataset are stored correctly in R as
factors. Copy and paste your output.

\subsection{(b)}\label{b-2}

Fit a model with the response (Y ) as credit card balance and X1= Income
and X2 = Student as the predictors. Call this model fit. Summarize your
output.

\subsection{(c)}\label{c-1}

Based on our results from part (b), write out the fitted model for

\subsubsection{i students}\label{i-students}

\subsubsection{ii non-students.}\label{ii-non-students.}

\subsection{(d)}\label{d-1}

Interpret the regression coefficient related to Income for both models.

\subsubsection{i students}\label{i-students-1}

\subsubsection{ii non-students.}\label{ii-non-students.-1}

\subsection{(e) (practice)}\label{e-practice}

Notice that our model says that regardless of student status, the effect
of Income on average Balance is the same. Do you think this is a
reasonable constraint on our model?

Construct some plots to back up your answer.

\#\#(f) (practice)

One way we could relax this assumption is by incorporating interaction
terms into our model. Specifically:

Yi = \$\beta\$0 + \$\beta\$1Xi1 + \$\beta\$2Xi2 + \$\beta\$3Xi3 +
\(\epsilon\)i,

where X1= Income, X2 = Student, and X3 = Income × Student. Fit a model
with an interaction term using the following code:

\begin{Shaded}
\begin{Highlighting}[]
\CommentTok{\# lm(Balance \textasciitilde{} Income + Student + Income:Student, data=Credit)}
\end{Highlighting}
\end{Shaded}

Based on this model, write out the fitted model for students and write
out the fitted model for non-students.

\subsection{(g) (practice)}\label{g-practice}

Interpret the regression coefficient related to Income for the fitted
models obtained in part (f).

\end{document}
